\section*{第 \ref{ch:b1:precipitation} 章——沉淀与溶解度}

\begin{solution}{exo:b1:precipitation:1}
\ce{BaSO4(s) <=> Ba^2+ + SO4^2-}，$K_s = [\ce{Ba^2+}][\ce{SO4^2-}] =
10^{-9.97}$。\ce{CaF2(s) <=> Ca^2+ + 2F-}，$K_s = [\ce{Ca^2+}][\ce{F-}]^2 =
10^{-9.84}$。\ce{Ag2CrO4(s) <=> 2Ag+ + CrO4^2-}，$K_s =
[\ce{Ag+}]^2[\ce{CrO4^2-}] = 10^{-11.95}$。\ce{Fe(OH)3(s) <=> Fe^3+ + 3OH-}，
$K_s = [\ce{Fe^3+}][\ce{OH-}]^3 = 10^{-38.55}$。对于 $K_s = \num{2.0e-13}$，
$\mathrm{p}K_s = 12.70$.
\end{solution}

\begin{solution}{exo:b1:precipitation:2}
$s = \sqrt{K_s} = 10^{-12.27/2} = \qty{7.3e-7}{mol/L}$，即
$\num{7.3e-7} \times 187.8 = \qty{1.4e-4}{g/L} = \qty{0.14}{mg/L}$。
\end{solution}

\begin{solution}{exo:b1:precipitation:3}
混合后（体积加倍），$[\ce{Ba^2+}] = \qty{1.0e-3}{mol/L}$，
$[\ce{SO4^2-}] = \qty{1.0e-4}{mol/L}$：$Q_s = \num{1.0e-7} > K_s =
\num{1.1e-10}$，硫酸钡沉淀。第二种情形中
$[\ce{SO4^2-}] = \num{1.0e-7}$，$Q_s = \num{1.0e-10} < K_s$，
无沉淀生成，但差距很小。
\end{solution}

\begin{solution}{exo:b1:precipitation:4}
在水中 $s = 10^{-9.97/2} = \qty{1.0e-5}{mol/L}$。在 \qty{0.010}{mol/L}
硫酸根溶液中（\cref{prop:b1:precipitation:common-ion}），$s = 10^{-9.97}/0.010 =
\qty{1.1e-8}{mol/L}$，约小一千倍（倍数为 970）。
\end{solution}

\begin{solution}{exo:b1:precipitation:5}
在水中 $s = (10^{-9.84}/4)^{1/3} = \qty{3.3e-4}{mol/L}$，即
\qty{26}{mg/L}。在 \qty{0.10}{mol/L} 氯化钙溶液中，
$[\ce{Ca^2+}] \approx 0.10$，$[\ce{F-}] = 2s$，因此
$s = \frac12\sqrt{10^{-9.84}/0.10} = \qty{1.9e-5}{mol/L}$。近似
$s \ll 0.10$ 成立，两者相差四个数量级。
\end{solution}

\begin{solution}{exo:b1:precipitation:6}
开始沉淀时：$\mathrm{pH} = 14.00 - \frac12(11.25 + \log 0.050) = 14.00 - 4.97 =
9.03$。沉淀率达 \qty{99}{\percent} 时，$[\ce{Mg^2+}] = \num{5.0e-4}$，
$\mathrm{pH} = 14.00 - \frac12(11.25 - 3.30) = 10.03$。在 pH 轴上，
\ce{Mg(OH)2} 在 9.03 以上存在；在其以下，镁完全以 \ce{Mg^2+}
形式存在，浓度为 \qty{0.050}{mol/L}。
\end{solution}

\begin{solution}{exo:b1:precipitation:7}
溴化银先沉淀，此时 $[\ce{Ag+}] = 10^{-12.27 + 2} =
10^{-10.27}$，氯化银在 $10^{-7.75}$ 时沉淀。当氯化银开始沉淀时，
$[\ce{Br-}] = 10^{-12.27}/10^{-7.75} = 10^{-4.52} =
\qty{3.0e-5}{mol/L}$，为溴离子的 \qty{0.30}{\percent}。
按通常的 \qty{0.1}{\percent} 判据，此分离不可接受：两者的
$\mathrm{p}K_s$ 仅相差 2.5。
\end{solution}

\begin{solution}{exo:b1:precipitation:8}
铬酸银在 $[\ce{Ag+}]^2 \times \num{5.0e-3} = 10^{-11.95}$ 时生成，
$[\ce{Ag+}] = \qty{1.5e-5}{mol/L}$。此时 $[\ce{Cl-}] = 10^{-9.75}/\num{1.5e-5}
= \qty{1.2e-5}{mol/L}$，占氯离子的 \num{2.4e-4}（\qty{0.024}{\percent}）。
红色出现时，氯离子实际上已
基本全部沉淀。
\end{solution}

\begin{solution}{exo:b1:precipitation:9}
$[\ce{CO3^2-}] = [\ce{HCO3-}]K_{a2}/[\ce{H3O+}] = [\ce{HCO3-}]\,10^{\mathrm{pH}
- \mathrm{p}K_{a2}}$，故 $K_s = [\ce{Ca^2+}][\ce{HCO3-}]\,10^{\mathrm{pH} -
\mathrm{p}K_{a2}}$，即所求关系式。取 $[\ce{Ca^2+}] =
[\ce{HCO3-}] = s$，则 $s^2 = 10^{-8.30 + 10.33 - 8.30} = 10^{-6.27}$，
$s = \qty{7.3e-4}{mol/L}$（\ce{CaCO3} 为 \qty{73}{mg/L}），若碳酸根离子不是碱，则为
$10^{-4.15} = \qty{7.1e-5}{mol/L}$。
\end{solution}

\begin{solution}{exo:b1:precipitation:10}
1.~加入 \ce{CO2(g) <=> CO2(aq)}：$K' = K K_H = 10^{-4.34 - 1.47} =
10^{-5.81}$。
2.~$[\ce{Ca^2+}] = s$，$[\ce{HCO3-}] = 2s$，故 $4s^3 = K'p$。在
\qty{0.010}{bar} 下：$s = (10^{-5.81} \times 0.010/4)^{1/3} =
\qty{1.6e-3}{mol/L}$，即 \qty{157}{mg/L}。在室外空气中：
$s = \qty{5.5e-4}{mol/L}$，即 \qty{55}{mg/L}。（验证：在土壤空气中
$[\ce{CO2(aq)}] = 10^{-3.47}$，$\mathrm{pH} = 6.37 + \log(3.14\times
10^{-3}/3.39 \times 10^{-4}) = 7.34$：碳酸氢根确实占优势。）
3.~在土壤中二氧化碳压强较高的条件下，水中已饱和了方解石，
这些水到达洞顶，那里压强较低，于是失去二氧化碳，
$Q > K'$，方解石便在水滴悬挂之处沉淀，
一圈圈地生成。
% ledger: air:CO2, dfg:CO2_g, dfg:CO2_ao (log KH computed in figdata/bachelor-1/aqueous-constants.py)
\end{solution}

\begin{solution}{exo:b1:precipitation:11}
1.~开始沉淀时 $\mathrm{pH} = 14.00 - \frac12(16.38 - 2.00) = 6.81$。完全
重新溶解时 $[\ce{Zn(OH)4^2-}] = 10^{-2} = 10^{-1.93}[\ce{OH-}]^2$，
$[\ce{OH-}] = 10^{-0.035}$，$\mathrm{pH} = 13.97$。氢氧化物存在于
pH 6.81 至 13.97 之间。
2.~$\mathrm{pH} = 14.00 - 14.45/4 = 10.39$，$s_{\min} = 2 \times
10^{-(16.38 + 1.93)/2} = \qty{1.4e-9}{mol/L}$（在此双物种模型中）。
3.~斜率为 $-2$ 的直线从 $(6.81, -2)$ 下降到极小值附近的
$(10.39, -8.9)$，再以斜率 $+2$ 上升到 $(13.97, -2)$；真实曲线
因其他羟基配合物的作用而圆滑化
（见\cref{sec:b1:precipitation:amphoteric}）。
\end{solution}

\begin{solution}{exo:b1:precipitation:12}
1.~$\ce{Al(OH)3(s) + OH- <=> Al(OH)4-}$，$K = 10^{-1.23}$；所有的
铝都溶解时，$10^{-3} \leqslant K[\ce{OH-}]$，即
$[\ce{OH-}] \geqslant 10^{-1.77}$，$\mathrm{pH} \geqslant 12.23$。
2.~$[\ce{Fe^3+}] = 10^{-38.55}/10^{-3 \times 1.77} = 10^{-33.2}$：铁
完全以 \ce{Fe(OH)3} 的形式存在；过滤将铁（固体）
与铝（滤液）分离。
3.~氢氧化镁不具两性：在过量的氢氧化物中，它仍为固体，
与氢氧化铁（III）一起，无法实现分离。
\end{solution}

\begin{solution}{pb:b1:precipitation:1}
\textbf{1.} \ce{Fe(OH)3(s) <=> Fe^3+ + 3OH-}, $K_s = [\ce{Fe^3+}][\ce{OH-}]^3$;
\ce{Zn(OH)2(s) <=> Zn^2+ + 2OH-}, $K_s = [\ce{Zn^2+}][\ce{OH-}]^2$;
\ce{Mg(OH)2(s) <=> Mg^2+ + 2OH-}，同理。
\textbf{2.} 这些常数的形式并不相同（$[\ce{OH-}]$ 的幂次为 3 或 2），
且三种离子的浓度也各不相同；
因此只能比较开始沉淀时的 pH。
\textbf{3.} 开始沉淀时 $c[\ce{OH-}]^n = K_s$，故 $n\log[\ce{OH-}] =
-\mathrm{p}K_s - \log c$，且 $\mathrm{pH} = \mathrm{p}K_e +
\log[\ce{OH-}] = \mathrm{p}K_e - \frac1n(\mathrm{p}K_s + \log c)$。
\textbf{4.} $14.00 - \frac13(38.55 - 3.00) = 2.15$.
\textbf{5.} $14.00 - \frac12(16.38 - 2.30) = 6.96$.
\textbf{6.} $14.00 - \frac12(11.25 - 1.70) = 9.22$.
\textbf{7.} 沉淀先后次序为：铁（III）（2.15）、锌（6.96）、镁（9.22）。
在 pH 为 2.0 时，对铁而言 $Q_s = 10^{-3} \times 10^{-36} = 10^{-39} < 10^{-38.55}$，
其余两种离子距沉淀更远，即无固体生成。
\textbf{8.} \qty{1.0e-6}{mol/L}.
\textbf{9.} $14.00 - \frac13(38.55 - 6.00) = 3.15$.
\textbf{10.} $\log[\ce{Fe^3+}] = -38.55 + 3(14.00 - \mathrm{pH}) = 3.45 -
3\,\mathrm{pH}$：斜率为 $-3$。
\textbf{11.} 直至氢氧化锌开始沉淀，pH 为 6.96。
\textbf{12.} pH 从 3.15 至 6.96。
\textbf{13.} \ce{Fe^3+ + OH- <=> FeOH^2+}, $\beta_1 =
[\ce{FeOH^2+}]/([\ce{Fe^3+}][\ce{OH-}]) = 10^{11.82}$.
\textbf{14.} 两者相等时 $[\ce{OH-}] = 10^{-11.82}$，即 pH 为 2.18：
\ce{Fe^3+} 在下方占优，\ce{FeOH^2+} 在上方占优。
\textbf{15.} $10^{11.82 + 3.15 - 14.00} = 10^{0.97} = 9.3$：pH 为 3.15 时，
溶解的铁为游离 \ce{Fe^3+} 的 10.3 倍，即 \qty{1.0e-5}{mol/L}：
仅能去除 \qty{99.0}{\percent}。第一次的估计并不可靠。
\textbf{16.} $[\ce{Fe}]_{\text{溶解}} = 10^{3.45 - 3\,\mathrm{pH}}\,
(1 + 10^{\mathrm{pH} - 2.18})$.
\textbf{17.} $\log[\ce{FeOH^2+}] = 11.82 + \log[\ce{Fe^3+}] + \mathrm{pH} -
14.00 = 1.27 - 2\,\mathrm{pH}$：斜率为 $-2$。
\textbf{18.} 由 $1.27 - 2\,\mathrm{pH} \approx -6.00$ 得 3.64；计入
\ce{Fe^3+} 项后，$10^{3.45 - 10.92}(1 + 10^{1.46}) =
\num{1.0e-6}$，此时 $\mathrm{pH} = 3.64$。
\textbf{19.} $[\ce{Zn^2+}] = \num{5.0e-5}$: $14.00 - \frac12(16.38 - 4.30) =
7.96$。氢氧化镁的开始沉淀 pH 为 9.22，此时尚不沉淀。
\textbf{20.} \ce{Zn(OH)2(s) + 2OH- <=> Zn(OH)4^2-}, $K = \beta_4 K_s =
10^{14.45 - 16.38} = 10^{-1.93}$.
\textbf{21.} $[\ce{OH-}]^2 = \num{5.0e-3}/10^{-1.93}$, $\mathrm{pH} =
13.81$。操作人员加碱不可过量：pH 高于 9.2 时，
镁会与锌一同沉淀，若远高于此，锌则
重新溶解。
\textbf{22.} $\num{1.0e-3} \times 1000 \times 0.999 \times 106.8 =
\qty{107}{g}$ 每立方米。
\textbf{23.} $\num{5.0e-3} \times 1000 \times 0.99 \times 99.4 =
\qty{492}{g}$ 每立方米。
\textbf{24.} 若一次升至 6.96 以上，两种氢氧化物将同时沉淀，
形成同一种污泥：锌随铁渣流失，铁渣
被锌污染。在两阈值之间停止，则可
单独去除铁。
\textbf{25.} \textbf{pH 从 3.64 至 6.96}：低于 3.64 时，仍有超过
\qty{0.1}{\percent} 的铁溶解，主要以 \ce{FeOH^2+} 形式存在；高于
6.96 时则生成氢氧化锌。忽略 \ce{FeOH^2+} 会使下界
定为 3.15。
\end{solution}
